\documentclass[12pt,a4paper]{article}

\usepackage[spanish]{babel}
\usepackage[utf8]{inputenc}
\usepackage[T1]{fontenc}
\usepackage{geometry}
\usepackage{amsmath}
\usepackage{amssymb}
\usepackage{booktabs}
\usepackage[ruled,vlined,linesnumbered]{algorithm2e}
\usepackage[hidelinks]{hyperref}
\usepackage{float}
\usepackage{graphicx}
\usepackage{tikz}
\usetikzlibrary{shapes.geometric, arrows.meta, positioning, fit, backgrounds, calc}
\newcommand{\maybeimage}[3]{%
  \IfFileExists{#1}{\includegraphics[#2]{#1}}%
  {\fbox{\parbox{0.7\textwidth}{\centering\vspace{1cm}[Figura pendiente: #3]\vspace{1cm}}}}%
}
\usepackage{caption}
\usepackage{parskip}
\usepackage{setspace}
\usepackage{titlesec}
\usepackage{xcolor}
\usepackage{array}
\usepackage{multirow}

\geometry{top=2.5cm, bottom=2.5cm, left=2.5cm, right=2.5cm}
\setlength{\parskip}{6pt}

\titlespacing*{\section}{0pt}{14pt}{6pt}
\titlespacing*{\subsection}{0pt}{10pt}{4pt}

% Custom colors for TikZ
\definecolor{colPop}{RGB}{52,101,164}
\definecolor{colSel}{RGB}{180,140,40}
\definecolor{colCross}{RGB}{39,126,77}
\definecolor{colMut}{RGB}{150,50,150}
\definecolor{colBL}{RGB}{200,50,50}
\definecolor{lightPop}{RGB}{225,235,255}
\definecolor{lightSel}{RGB}{255,245,210}
\definecolor{lightCross}{RGB}{220,245,230}
\definecolor{lightMut}{RGB}{245,230,255}
\definecolor{lightBL}{RGB}{255,220,220}

\newcommand{\placeholder}[0]{\textit{---}}

\SetKwInput{KwIn}{Entrada}
\SetKwInput{KwOut}{Salida}
\SetKw{KwReturn}{devolver}
\SetKwFor{Para}{para}{hacer}{fin para}
\SetKwFor{Mientras}{mientras}{hacer}{fin mientras}
\SetKwIF{Si}{SiNo}{OSiNo}{si}{entonces}{si no si}{si no}{fin si}
\SetKwFunction{fitness}{fitness}
\SetKwFunction{fix}{fix}
\SetKwFunction{createSolution}{crearSolución}
\SetKwFunction{torneo}{torneo}
\SetKwFunction{cruzar}{cruzar}
\SetKwFunction{mutar}{mutar}
\SetKwFunction{BL}{BL}

%-----------------------------------------------------------------------
\begin{document}

%-----------------------------------------------------------------------
% PORTADA
%-----------------------------------------------------------------------
\begin{titlepage}
  \centering
  \vspace*{2cm}
  {\Large \textbf{Universidad de Granada}\\[0.3cm]
          \textbf{ETSIIT Informática y Telecomunicación}\\[0.3cm]
          Dpto.\ de Ciencias de la Computación e Inteligencia Artificial}

  \vspace{2.5cm}
  {\LARGE \textbf{Metaheurísticas}}\\[0.5cm]
  {\Large Práctica 2}\\[0.3cm]
  {\large Técnicas de Búsqueda basadas en Poblaciones\\
           para el Problema de Portfolio}

  \vspace{1.5cm}
  {\large Curso 2025--2026}

  \vfill
  \begin{tabular}{rl}
    \textbf{Nombre:}  & Jose Angel Carretero Montes \\[4pt]
    \textbf{DNI:}     & 25628892S \\[4pt]
    \textbf{Email:}   & \texttt{joseangel1@correo.ugr.es} \\[4pt]
    \textbf{Grupo:}   & Grupo 3 \\[4pt]
    \textbf{Horario:} & Martes 15:30--17:30 \\
  \end{tabular}
  \vspace{1.5cm}
\end{titlepage}

%-----------------------------------------------------------------------
% ÍNDICE
%-----------------------------------------------------------------------
\tableofcontents
\newpage

%-----------------------------------------------------------------------
\section{Descripción del Problema}
%-----------------------------------------------------------------------

El \textbf{Problema de Portfolio} consiste en determinar la asignación óptima de capital
entre $N$ activos financieros con el objetivo de maximizar el retorno esperado
minimizando simultáneamente el riesgo.

\subsection*{Representación de los datos}

Para cada activo $i \in \{1,\ldots,N\}$ y día $t \in \{1,\ldots,T\}$ se dispone
del retorno diario $\mu_{ti}$, calculado como el ratio de mejora relativa del valor:
\[
  \mu_{ti} = \frac{V_{ti} - V_{(t+1)i}}{V_{ti}}
\]
Los datos empleados cubren el periodo 2015--2024 para entrenamiento y el año 2025
para evaluación de generalización.

\subsection*{Función objetivo}

Dado un vector de pesos $w = (w_1,\ldots,w_N)$, la calidad de la cartera se mide como:
\[
  \text{fitness}(w) = \text{Beneficio}(w) - \lambda \cdot \text{Riesgo}(w)
\]
donde:
\[
  \text{Beneficio}(w) = \sum_{i=1}^{N} w_i \cdot \sum_{t=1}^{T}\log(1+\mu_{ti})
  \qquad
  \text{Riesgo}(w) = \sqrt{\sum_{i=1}^{N}\sum_{j=1}^{N} w_i \cdot \Sigma_{ij} \cdot w_j}
\]
siendo $\Sigma_{ij}$ la covarianza entre activos $i$ y $j$, y $\lambda = 500$ el factor de
penalización del riesgo.

\subsection*{Restricciones}

\[
  \sum_{i=1}^{N} w_i = 1 \qquad
  w_i = 0 \;\vee\; lo \le w_i \le hi \quad \forall i \in 1\ldots N
\]

Los rangos $[lo, hi]$ varían según el conjunto de datos (véase Sección~\ref{sec:experimentos}).

%-----------------------------------------------------------------------
\section{Aplicación de los Algoritmos al Problema}
%-----------------------------------------------------------------------

\subsection{Representación de soluciones}

Una solución queda representada como un vector $w \in \mathbb{R}^N$
donde cada componente $w_i$ indica la fracción del capital invertida en el activo $i$.
La restricción de cardinalidad implica que $w_i$ puede ser exactamente cero
(no se invierte en esa empresa) o estar en el intervalo $[lo, hi]$.

\subsection{Generación aleatoria de soluciones}

La generación de la solución inicial es la misma que en la Práctica 1: se elige un
subconjunto aleatorio de empresas activas, se les asigna un peso aleatorio en $[lo, hi]$,
se normaliza y se aplica el operador de reparación.

\subsection{Operador de reparación (\texttt{fix})}

El operador de reparación es el mismo que el empleado en la Práctica 1 (tres fases:
normalización, recorte de rango, redistribución iterativa).
Es especialmente importante en los algoritmos genéticos ya que los operadores de
cruce pueden generar soluciones que violen las restricciones.

\subsection{Operador de selección: torneo binario}

En ambos esquemas (generacional y estacionario) se usa el \textbf{torneo binario}:
se eligen dos individuos al azar de la población y se devuelve el de mayor fitness.

\begin{algorithm}[H]
\caption{Torneo binario}
\KwIn{población $P$}
\KwOut{individuo ganador}
$i_1 \leftarrow$ índice aleatorio en $[0, |P|-1]$\;
$i_2 \leftarrow$ índice aleatorio en $[0, |P|-1]$\;
\Si{$P[i_1].\text{fitness} \geq P[i_2].\text{fitness}$}{
  \KwReturn{$P[i_1]$}
}
\KwReturn{$P[i_2]$}
\end{algorithm}

\subsection{Operadores de cruce}

Se emplean dos operadores para representación real.

\subsubsection*{Cruce Aritmético}

Para cada posición $i$ se genera un $r_i \in [0,1]$ independiente:
\[
  \text{Off}^1_i = r_i \cdot S^1_i + (1-r_i) \cdot S^2_i
  \qquad
  \text{Off}^2_i = r_i \cdot S^2_i + (1-r_i) \cdot S^1_i
\]

\subsubsection*{Cruce BLX-$\alpha$}

Para cada posición $i$, siendo $I_i = |S^1_i - S^2_i|$:
\[
  hmin_i = \min(S^1_i, S^2_i) - \alpha \cdot I_i
  \qquad
  hmax_i = \max(S^1_i, S^2_i) + \alpha \cdot I_i
\]
Ambas soluciones hijas se generan de igual forma: cada gen del hijo se muestrea
uniformemente en $[\max(0, hmin_i), \min(hi, hmax_i)]$, con $\alpha = 0{,}3$.

Tras ambos cruces se aplica \textbf{fix} a los hijos generados.

\subsection{Operador de mutación}

\subsubsection*{Para Algoritmos Genéticos (AG)}
Con probabilidad $p_m = 0{,}1$ \textbf{por individuo}, se selecciona un gen aleatorio
y se reemplaza por un valor uniforme en $[lo, hi]$, seguido de \textbf{fix}.
Solo cuenta como evaluación si efectivamente se muta.

\subsubsection*{Para Algoritmos Meméticos (AM)}
Con probabilidad $p_m = 0{,}1/N$ \textbf{por gen}, cada gen se reemplaza independientemente
por un valor uniforme en $[lo, hi]$, seguido de \textbf{fix}.
Solo cuenta como evaluación si efectivamente se muta algún gen.

La mutación por gen en AMs favorece una exploración más uniforme del espacio,
mientras que la mutación por individuo en AGs es más selectiva.

%-----------------------------------------------------------------------
\section{Algoritmos Genéticos y Meméticos}
%-----------------------------------------------------------------------

\subsection{Algoritmo Genético Generacional (AGG)}

El AGG mantiene una población completa de $N_p = 100$ individuos que se reemplaza
generación a generación. Incorpora elitismo para no perder la mejor solución encontrada.

\subsubsection*{Parámetros}
\begin{itemize}
  \item Tamaño de población: $N_p = 100$
  \item Probabilidad de cruce: $p_c = 0{,}8$
  \item Probabilidad de mutación: $p_m = 0{,}1$ (por individuo)
  \item Criterio de parada: $10\,000$ evaluaciones de la función objetivo
\end{itemize}

\subsubsection*{Esquema de evolución y reemplazamiento}

\begin{algorithm}[H]
\caption{AGG (esquema generacional con elitismo)}
\KwIn{problema, maxEvals}
\KwOut{mejor solución encontrada}
Inicializar población $P$ con $N_p$ individuos aleatorios\;
$evals \leftarrow N_p$\;
\Mientras{$evals < \text{maxEvals}$}{
  $P' \leftarrow \emptyset$\;
  \Para{$i \leftarrow 1$ \KwTo $N_p/2$ \textbf{y} $evals < \text{maxEvals}$}{
    $p_1 \leftarrow$ \torneo{$P$};\quad $p_2 \leftarrow$ \torneo{$P$}\;
    \Si{$\text{rand}() \leq 0{,}8$}{
      $(h_1, h_2) \leftarrow$ \cruzar{$p_1, p_2$}; \quad $evals \leftarrow evals + 2$\;
    }\SiNo{
      $h_1 \leftarrow p_1$;\quad $h_2 \leftarrow p_2$\;
    }
    \mutar{$h_1$}; \quad \mutar{$h_2$}\;
    $P' \leftarrow P' \cup \{h_1, h_2\}$\;
  }
  \tcp{Elitismo: conservar el mejor padre si la nueva generación es peor}
  \Si{$\max(P').\text{fitness} < \max(P).\text{fitness}$}{
    $\min(P') \leftarrow \max(P)$\;
  }
  $P \leftarrow P'$\;
}
\KwReturn{$\max(P)$}
\end{algorithm}

\noindent Se implementan dos variantes: \textbf{AGG-Arit} (cruce aritmético) y
\textbf{AGG-BLX} (cruce BLX-$\alpha$), siendo esa la única diferencia entre ellas.

\begin{figure}[H]
\centering
\begin{tikzpicture}[
  node distance=0.4cm and 0.85cm,
  blk/.style={rectangle, rounded corners=4pt, minimum width=1.9cm, minimum height=0.85cm,
              text centered, text width=1.9cm, font=\scriptsize\bfseries, thick},
  arr/.style={-{Stealth[length=2.5mm]}, thick}
]
\node[blk, fill=lightPop, draw=colPop] (P) {Población $P(t)$ \\ $N_p=100$};
\node[blk, fill=lightSel, draw=colSel, right=of P] (sel) {Selección \\ Torneo ($\times 2$)};
\node[blk, fill=lightCross, draw=colCross, right=of sel] (cru) {Cruce \\ $p_c=0.8$};
\node[blk, fill=lightMut, draw=colMut, right=of cru] (mut) {Mutación \\ $p_m=0.1$/ind};
\node[blk, fill=lightPop, draw=colPop, right=of mut] (Pp) {Nueva $P'$ \\ $N_p=100$};

\node[blk, fill=lightBL, draw=colBL, below=1.5cm of cru] (elit) {Elitismo \\ mejor$(P) \to$ peor$(P')$};

\draw[arr] (P) -- (sel);
\draw[arr] (sel) -- (cru);
\draw[arr] (cru) -- (mut);
\draw[arr] (mut) -- (Pp);

\draw[arr, dashed] (P.south) |- (elit.west);
\draw[arr, dashed] (elit.east) -| (Pp.south);

\draw[arr, rounded corners=5pt] (Pp.east) -- ++(0.6,0) -- ++(0,1.2)
      -- node[above, font=\tiny]{$t \to t+1$} ($ (P.west)+(-0.6,1.2) $)
      -- ($ (P.west)+(-0.6,0) $) -- (P.west);
\end{tikzpicture}
\caption{Ciclo evolutivo del AGG. Las líneas discontinuas muestran el flujo de elitismo.}
\label{fig:agg_cycle}
\end{figure}

\subsection{Algoritmo Genético Estacionario (AGE)}

El AGE genera solo dos descendientes por iteración y los inserta en la población
reemplazando a los peores individuos si mejoran. La población apenas varía
en cada paso, lo que lo hace más conservador.

\subsubsection*{Parámetros}
\begin{itemize}
  \item Tamaño de población: $N_p = 100$
  \item Probabilidad de cruce: $p_c = 1{,}0$ (siempre se cruzan)
  \item Probabilidad de mutación: $p_m = 0{,}1$ (por individuo)
  \item Criterio de parada: $10\,000$ evaluaciones de la función objetivo
\end{itemize}

\subsubsection*{Esquema de evolución y reemplazamiento}

\begin{algorithm}[H]
\caption{AGE (esquema estacionario)}
\KwIn{problema, maxEvals}
\KwOut{mejor solución encontrada}
Inicializar población $P$ con $N_p$ individuos aleatorios\;
$evals \leftarrow N_p$\;
\Mientras{$evals < \text{maxEvals}$}{
  $p_1 \leftarrow$ \torneo{$P$};\quad $p_2 \leftarrow$ \torneo{$P$}\;
  $(h_1, h_2) \leftarrow$ \cruzar{$p_1, p_2$}; \quad $evals \leftarrow evals + 2$\;
  \mutar{$h_1$}; \quad \mutar{$h_2$}\;
  \tcp{Reemplazamiento competitivo: cada hijo desafía al peor}
  \Si{$h_1.\text{fitness} > \min(P).\text{fitness}$}{
    $\min(P) \leftarrow h_1$\;
  }
  \Si{$h_2.\text{fitness} > \min(P).\text{fitness}$}{
    $\min(P) \leftarrow h_2$\;
  }
}
\KwReturn{$\max(P)$}
\end{algorithm}

\noindent Se implementan dos variantes: \textbf{AGE-Arit} y \textbf{AGE-BLX}.

\begin{figure}[H]
\centering
\begin{tikzpicture}[
  node distance=0.4cm and 0.85cm,
  blk/.style={rectangle, rounded corners=4pt, minimum width=1.9cm, minimum height=0.85cm,
              text centered, text width=1.9cm, font=\scriptsize\bfseries, thick},
  arr/.style={-{Stealth[length=2.5mm]}, thick}
]
\node[blk, fill=lightPop, draw=colPop] (P) {Población $P$ \\ $N_p=100$};
\node[blk, fill=lightSel, draw=colSel, right=of P] (sel) {Selección \\ Torneo ($\times 2$)};
\node[blk, fill=lightCross, draw=colCross, right=of sel] (cru) {Cruce \\ $p_c=1{,}0$};
\node[blk, fill=lightMut, draw=colMut, right=of cru] (mut) {Mutación \\ $p_m=0.1$/ind};
\node[blk, fill=lightPop, draw=colPop, right=of mut] (hijos) {2 hijos \\ $(h_1,h_2)$};

\node[blk, fill=lightBL, draw=colBL, below=1.5cm of cru] (rep)
      {Reemplazo \\ $h_i >$ peor$(P)$?};

\draw[arr] (P) -- (sel);
\draw[arr] (sel) -- (cru);
\draw[arr] (cru) -- (mut);
\draw[arr] (mut) -- (hijos);

% hijos compiten en reemplazo
\draw[arr] (hijos.south) |- (rep.east);

% P consultado para encontrar al peor (dashed)
\draw[arr, dashed] (P.south) |- (rep.west);

% reemplazo actualiza P en sitio: lazo bajo el diagrama (dashed)
\draw[arr, dashed, rounded corners=5pt] (rep.south) -- ++(0,-0.45) -| (P.south);
\end{tikzpicture}
\caption{Ciclo evolutivo del AGE. Las líneas discontinuas muestran la consulta del peor
individuo de $P$ y la actualización en sitio: solo 0--2 miembros cambian por iteración.}
\label{fig:age_cycle}
\end{figure}

\subsection{Algoritmos Meméticos (AM): Integración de Búsqueda Local}

Los Algoritmos Meméticos combinan la exploración global de los AGs con la explotación
local de la Búsqueda Local (BL). La BL se aplica \textbf{cada 10 generaciones}
a subconjuntos de la población según tres estrategias de selección.

Los AMs se basan en el AGG que mejores resultados ha proporcionado, que según los
experimentos es \textbf{AGG-BLX} (véase Sección~\ref{sec:experimentos}).

\subsubsection*{Búsqueda Local integrada}

La BL empleada es la \textbf{primero el mejor} (first improvement) con el operador
de movimiento del portfolio: transferir una fracción $p=0{,}15$ del peso
de un activo origen a un activo destino. Parámetro: máximo 100 evaluaciones por llamada.
Esta es la misma BL de la Práctica 1 pero con fracción de movimiento reducida del 40\% al 15\%.
El contador total de evaluaciones incluye las realizadas durante la BL.

\subsubsection*{Parámetros comunes de AMs}
\begin{itemize}
  \item Tamaño de población: $N_p = 100$
  \item Esquema base: AGG-BLX (cruce BLX-$\alpha$, $p_c = 0{,}8$)
  \item Probabilidad de mutación: $p_m = 0{,}1/N$ (por gen)
  \item Aplicación de BL: cada $10$ generaciones
  \item Presupuesto por llamada BL: $\min(100,\; \text{maxEvals} - evals)$ evaluaciones
  \item Criterio de parada: $10\,000$ evaluaciones (incluyendo las de la BL)
\end{itemize}

\subsubsection*{AM-All: Búsqueda Local a todos los individuos}

\begin{algorithm}[H]
\caption{AM-All: Aplicar BL a toda la población}
\KwIn{población $P$, problem, maxevals\_local}
\Mientras{$evals < \text{maxEvals}$}{
  \tcp{Generación (esquema generacional AGG-BLX)}
  \tcp{...}
  \Si{$generaciones \bmod 10 = 0$}{
    \Para{$i \leftarrow 1$ \KwTo $|P|$ \textbf{y} $evals < \text{maxEvals}$}{
      $budget \leftarrow \min(maxevals\_local,\; \text{maxEvals} - evals)$\;
      $\text{result} \leftarrow$ BL($\text{problem}$, $P[i]$, $budget$)\;
      $P[i] \leftarrow \text{result}$\;
      $evals \leftarrow evals + \text{result.evals}$\;
    }
  }
}
\end{algorithm}

AM-All aplica BL a tantos individuos como el presupuesto permita (típicamente $\approx 90$
de los 100 en la primera fase). Concentra toda la explotación local en una etapa
con poca evolución genética previa.

\subsubsection*{AM-Rand: Búsqueda Local aleatoria (10\% de población)}

\begin{algorithm}[H]
\caption{AM-Rand: Aplicar BL a 10\% aleatorio}
\KwIn{población $P$, problem, maxevals\_local}
\Mientras{$evals < \text{maxEvals}$}{
  \tcp{Generación (esquema generacional AGG-BLX)}
  \tcp{...}
  \Si{$generaciones \bmod 10 = 0$}{
    \Para{$i \leftarrow 1$ \KwTo $|P|$ \textbf{y} $evals < \text{maxEvals}$}{
      \Si{$\text{rand}() \leq 0{,}1$}{
        $budget \leftarrow \min(maxevals\_local,\; \text{maxEvals} - evals)$\;
        $\text{result} \leftarrow$ BL($\text{problem}$, $P[i]$, $budget$)\;
        $P[i] \leftarrow \text{result}$\;
        $evals \leftarrow evals + \text{result.evals}$\;
      }
    }
  }
}
\end{algorithm}

AM-Rand selecciona aleatoriamente aproximadamente el 10\% de la población
(esperanza matemática: $\approx 10$ individuos de 100) para aplicar BL, balanceando
exploración genética con explotación local.

\subsubsection*{AM-Best: Búsqueda Local a los $0{,}1 \cdot N_p$ mejores}

\begin{algorithm}[H]
\caption{AM-Best: Aplicar BL a élite}
\KwIn{población $P$, problem, maxevals\_local}
\Mientras{$evals < \text{maxEvals}$}{
  \tcp{Generación (esquema generacional AGG-BLX)}
  \tcp{...}
  \Si{$generaciones \bmod 10 = 0$}{
    \textbf{sort}($P$, descendente por fitness)\;
    \Para{$i \leftarrow 1$ \KwTo $\min(0{,}1 \cdot |P|,\; |P|)$ \textbf{y} $evals < \text{maxEvals}$}{
      $budget \leftarrow \min(maxevals\_local,\; \text{maxEvals} - evals)$\;
      $\text{result} \leftarrow$ BL($\text{problem}$, $P[i]$, $budget$)\;
      $P[i] \leftarrow \text{result}$\;
      $evals \leftarrow evals + \text{result.evals}$\;
    }
  }
}
\end{algorithm}

AM-Best aplica BL solamente a los $0{,}1 \cdot N_p = 10$ mejores individuos, focalizando
el presupuesto de explotación en las soluciones más prometedoras.

\begin{figure}[H]
\centering
\begin{tikzpicture}
\tikzset{
  cell/.style={rectangle, minimum width=0.28cm, minimum height=0.55cm, draw=black!30, thick},
  filled/.style={fill=colCross, draw=colCross, thick},
  empty/.style={fill=white, draw=black!20}
}
% Offset: labels at x=0 (right-anchor), cells start at x=2.0
\def\cellstart{2.0}
\def\cellstep{0.32}

% Labels (left edge at x=0, text extends right)
\node[right, font=\small\bfseries] at (0, 1.5) {AM-All:};
\node[right, font=\small\bfseries] at (0, 0.5) {AM-Rand:};
\node[right, font=\small\bfseries] at (0, -0.5) {AM-Best:};

% AM-All: all cells filled
\foreach \x in {0,1,...,19} {
  \node[cell, filled] at (\cellstart+\x*\cellstep, 1.5) {};
}

% AM-Rand: ~2 cells filled (10% of 20)
\foreach \x in {0,1,...,19} {
  \node[cell, empty] at (\cellstart+\x*\cellstep, 0.5) {};
}
\node[cell, filled] at (\cellstart+2*\cellstep, 0.5) {};
\node[cell, filled] at (\cellstart+15*\cellstep, 0.5) {};

% AM-Best: first 2 cells filled (10% of 20), sorted
\foreach \x in {0,1,...,19} {
  \node[cell, empty] at (\cellstart+\x*\cellstep, -0.5) {};
}
\node[cell, filled] at (\cellstart+0*\cellstep, -0.5) {};
\node[cell, filled] at (\cellstart+1*\cellstep, -0.5) {};
\node[font=\tiny, right] at (\cellstart+19.7*\cellstep, -0.5) {$\leftarrow$ mejores $\cdots$ peores};

% Legend
\node[font=\tiny, fill=colCross, inner sep=2pt] at (\cellstart+21*\cellstep, 1.5) {Aplicar BL};
\node[font=\tiny, fill=white, draw=black!20, inner sep=2pt] at (\cellstart+21*\cellstep, 1.0) {Sin BL};

% Note about scale
\node[font=\scriptsize] at (\cellstart+10*\cellstep, -1.4) {(Representación a escala 1:5 de población N$_p$=100)};
\end{tikzpicture}
\caption{Comparación de estrategias de selección en AMs. Cada célula representa 5 individuos.}
\label{fig:am_strategies}
\end{figure}

%-----------------------------------------------------------------------
\section{Estructura del Código y Manual de Usuario}
%-----------------------------------------------------------------------

\subsection{Estructura de directorios}

\begin{verbatim}
portfolio/
  main.cpp                  # Experimentos de 50 semillas
  CMakeLists.txt            # Sistema de construcción
  data/
    ibex_35.csv             # 30 empresas, 2015-2025
    syp_100.csv             # 97 empresas, 2015-2025
    syp_500.csv             # 457 empresas, 2015-2025
  common/
    mh.h, problem.h, random.hpp, util.h
  inc/
    portfolioProblem.h, randomSearch.h, greedy.h
    localsearch.h, genetic.h
  src/
    portfolioProblem.cpp, randomSearch.cpp, greedy.cpp
    localsearch.cpp, genetic.cpp
\end{verbatim}

\subsection{Compilación y ejecución}

\begin{verbatim}
cmake .
make
./main [max_evals] [n_seeds]
\end{verbatim}

Ejemplos: \texttt{./main} (10000 evals, 50 semillas),
\texttt{./main 5000 30} (5000 evals, 30 semillas).

%-----------------------------------------------------------------------
\section{Experimentos y Resultados}
\label{sec:experimentos}
%-----------------------------------------------------------------------

\subsection{Casos del problema}

\begin{center}
\begin{tabular}{lrrrl}
\toprule
\textbf{Fichero} & \textbf{Empresas ($N$)} & \textbf{Días train} & \textbf{Días test} & \textbf{Rango $[lo, hi]$} \\
\midrule
\texttt{ibex\_35.csv}  & 30  & 2559 & 254 & $[0.005,\; 0.08]$ \\
\texttt{syp\_100.csv}  & 97  & 2516 & 249 & $[0.005,\; 0.05]$ \\
\texttt{syp\_500.csv}  & 457 & 2516 & 249 & $[0.005,\; 0.02]$ \\
\bottomrule
\end{tabular}
\end{center}

Parámetros comunes: $\lambda = 500$, MAX\_EVALS $= 10\,000$, 50 semillas ($1$--$50$).

\subsection{Resultados}

\begin{table}[H]
\centering
\footnotesize
\caption{Resultados IBEX 35 --- $N=30$, $[lo,hi]=[0.005, 0.08]$ (medias, 50 semillas)}
\begin{tabular}{lrrrrr}
\toprule
\textbf{Algoritmo} & \textbf{Fit.Train} & \textbf{Fit.2025} & \textbf{Ben.2025} & \textbf{Evals} & \textbf{Tiempo (s)} \\
\midrule
Greedy   & $-4.6495$ & $-3.3868$ & $0.1152$ & $1.0$ & $0.000$ \\
Random   & $-4.9411$ & $-3.3060$ & $0.0360$ & $10000.0$ & $0.024$ \\
BL       & $-4.4912$ & $-3.0850$ & $0.0293$ & $2115.8$ & $0.001$ \\
AGG-Arit & $-4.7065$ & $-3.1564$ & $0.0408$ & $10000.6$ & $0.006$ \\
AGG-BLX  & $-4.5586$ & $-3.1134$ & $0.0377$ & $10000.6$ & $0.008$ \\
AGE-Arit & $-4.6010$ & $-3.1188$ & $0.0347$ & $10000.7$ & $0.008$ \\
AGE-BLX  & $-4.5205$ & $-3.1024$ & $0.0342$ & $10000.7$ & $0.008$ \\
AM-All   & $-4.8652$ & $-3.2068$ & $0.0473$ & $10000.0$ & $0.025$ \\
AM-Rand  & $-4.5161$ & $-3.1009$ & $0.0381$ & $10000.2$ & $0.016$ \\
AM-Best  & $-4.5026$ & $-3.0980$ & $0.0373$ & $10000.4$ & $0.016$ \\
\bottomrule
\end{tabular}
\end{table}

\begin{table}[H]
\centering
\footnotesize
\caption{Resultados S\&P 100 --- $N=97$, $[lo,hi]=[0.005, 0.05]$ (medias, 50 semillas)}
\begin{tabular}{lrrrrr}
\toprule
\textbf{Algoritmo} & \textbf{Fit.Train} & \textbf{Fit.2025} & \textbf{Ben.2025} & \textbf{Evals} & \textbf{Tiempo (s)} \\
\midrule
Greedy   & $-3.6635$ & $-5.0204$ & $0.0099$ & $1.0$ & $0.000$ \\
Random   & $-3.3904$ & $-4.1131$ & $0.0799$ & $10000.0$ & $0.064$ \\
BL       & $-2.7911$ & $-3.7565$ & $0.1088$ & $10000.0$ & $0.039$ \\
AGG-Arit & $-3.0938$ & $-3.8300$ & $0.0868$ & $10000.5$ & $0.030$ \\
AGG-BLX  & $-3.0739$ & $-3.8198$ & $0.0851$ & $10000.5$ & $0.037$ \\
AGE-Arit & $-3.0304$ & $-3.8065$ & $0.0888$ & $10000.7$ & $0.029$ \\
AGE-BLX  & $-2.9684$ & $-3.7770$ & $0.0903$ & $10000.5$ & $0.034$ \\
AM-All   & $-3.4100$ & $-4.0801$ & $0.0743$ & $10000.0$ & $0.221$ \\
AM-Rand  & $-3.1122$ & $-3.8244$ & $0.0809$ & $10000.2$ & $0.140$ \\
AM-Best  & $-3.0789$ & $-3.8091$ & $0.0857$ & $10000.6$ & $0.133$ \\
\bottomrule
\end{tabular}
\end{table}

\begin{table}[H]
\centering
\footnotesize
\caption{Resultados S\&P 500 --- $N=457$, $[lo,hi]=[0.005, 0.02]$ (medias, 50 semillas)}
\begin{tabular}{lrrrrr}
\toprule
\textbf{Algoritmo} & \textbf{Fit.Train} & \textbf{Fit.2025} & \textbf{Ben.2025} & \textbf{Evals} & \textbf{Tiempo (s)} \\
\midrule
Greedy   & $-3.4384$ & $-5.2046$ & $-0.0628$ & $1.0$ & $0.000$ \\
Random   & $-3.6419$ & $-4.4297$ & $0.0123$ & $10000.0$ & $1.734$ \\
BL       & $-2.6282$ & $-3.9162$ & $0.0232$ & $10000.0$ & $1.168$ \\
AGG-Arit & $-2.9723$ & $-3.9803$ & $0.0170$ & $10000.5$ & $0.656$ \\
AGG-BLX  & $-2.8710$ & $-3.9238$ & $0.0192$ & $10000.8$ & $0.673$ \\
AGE-Arit & $-2.9693$ & $-3.9774$ & $0.0163$ & $10000.6$ & $0.654$ \\
AGE-BLX  & $-2.8756$ & $-3.9346$ & $0.0184$ & $10000.7$ & $0.663$ \\
AM-All   & $-3.4839$ & $-4.2879$ & $0.0155$ & $10000.0$ & $2.111$ \\
AM-Rand  & $-2.9504$ & $-3.9267$ & $0.0223$ & $10000.2$ & $1.542$ \\
AM-Best  & $-2.9334$ & $-3.9337$ & $0.0183$ & $10000.6$ & $1.326$ \\
\bottomrule
\end{tabular}
\end{table}

\subsection{Boxplots}

Las Figuras~\ref{fig:box_ibex}--\ref{fig:box_syp500} muestran la distribución del
fitness de entrenamiento en las 50 ejecuciones para cada conjunto de datos.

\begin{figure}[H]
  \centering
  \maybeimage{results_ibex_35.png}{width=0.85\textwidth}{results\_ibex\_35.png}
  \caption{Distribución del fitness (2015--2024) en IBEX 35.}
  \label{fig:box_ibex}
\end{figure}

\begin{figure}[H]
  \centering
  \maybeimage{results_syp_100.png}{width=0.85\textwidth}{results\_syp\_100.png}
  \caption{Distribución del fitness (2015--2024) en S\&P 100.}
  \label{fig:box_syp100}
\end{figure}

\begin{figure}[H]
  \centering
  \maybeimage{results_syp_500.png}{width=0.85\textwidth}{results\_syp\_500.png}
  \caption{Distribución del fitness (2015--2024) en S\&P 500.}
  \label{fig:box_syp500}
\end{figure}

%-----------------------------------------------------------------------
\section{Análisis de Resultados}
%-----------------------------------------------------------------------

\subsection{Impacto del operador de cruce: AGG-Arit vs AGG-BLX}

El cruce BLX-$\alpha$ con $\alpha=0.3$ genera descendientes que pueden extenderse
\textbf{más allá} del intervalo entre los padres, permitiendo explorar nuevas regiones
del espacio. El cruce aritmético, al ser una combinación lineal, permanece siempre
dentro del intervalo convexo de los padres.

\begin{center}
\begin{tabular}{lccc}
\toprule
\textbf{Dataset} & \textbf{AGG-Arit} & \textbf{AGG-BLX} & \textbf{Mejora BLX} \\
\midrule
ibex\_35 (N=30) & $-4.7065$ & $-4.5586$ & $0.148$ \\
syp\_100 (N=97) & $-3.0938$ & $-3.0739$ & $0.020$ \\
syp\_500 (N=457) & $-2.9723$ & $-2.8710$ & $0.101$ \\
\bottomrule
\end{tabular}
\end{center}

\textbf{Observación clave}: BLX supera a Aritmético en todos los datasets. La ventaja es
más pronunciada en problemas pequeños (ibex\_35) y grandes (syp\_500). En espacios
pequeños, BLX escapa más fácilmente de óptimos locales; en espacios grandes, la libertad
exploratoria de BLX es crítica para encontrar soluciones competitivas.

El cruce aritmético es más explotativo (refina dentro de la región muestreada),
mientras que BLX es más explorativo (descubre nuevas regiones del espacio).
Estos resultados justifican que los AMs usen \textbf{AGG-BLX} como base.

\subsection{Impacto del operador de cruce: AGE-Arit vs AGE-BLX}

El mismo patrón se observa en el esquema estacionario:

\begin{center}
\begin{tabular}{lccc}
\toprule
\textbf{Dataset} & \textbf{AGE-Arit} & \textbf{AGE-BLX} & \textbf{Mejora BLX} \\
\midrule
ibex\_35 & $-4.6010$ & $-4.5205$ & $0.081$ \\
syp\_100 & $-3.0304$ & $-2.9684$ & $0.062$ \\
syp\_500 & $-2.9693$ & $-2.8756$ & $0.094$ \\
\bottomrule
\end{tabular}
\end{center}

BLX mejora consistentemente AGE en todos los datasets. Esto confirma que el operador de cruce
es un factor crítico independiente del esquema de reemplazamiento.

\subsection{Comparación de esquemas: AGG vs AGE}

Usando BLX para una comparación justa:

\begin{center}
\begin{tabular}{lcc|c}
\toprule
\textbf{Dataset} & \textbf{AGG-BLX} & \textbf{AGE-BLX} & \textbf{Ganador} \\
\midrule
ibex\_35 & $-4.5586$ & $-4.5205$ & AGE ($+0.038$) \\
syp\_100 & $-3.0739$ & $-2.9684$ & AGE ($+0.106$) \\
syp\_500 & $-2.8710$ & $-2.8756$ & AGG ($+0.005$) \\
\bottomrule
\end{tabular}
\end{center}

\textbf{Conclusión}: AGE es superior a AGG en problemas de tamaño pequeño a medio.
En problemas grandes (syp\_500), son prácticamente equivalentes (diferencia de solo $0.005$).

\textbf{Razón algorítmica}: El esquema generacional AGG reemplaza \textbf{toda} la población
cada generación, causando \textit{drift genético}: pérdida de soluciones buenas que no fueron
seleccionadas en esa generación, pese al elitismo.
AGE, por el contrario, mantiene 98 individuos intactos y solo reemplaza los 2 peores,
conservando naturalmente una élite implícita. Con un presupuesto limitado (10000 evals),
la estrategia conservadora de AGE es superior.

En problemas grandes, el drift en AGG queda compensado por la exploración masiva que
permite (50 parejas nuevas por generación vs.\ 2 en AGE).

\subsection{Algoritmos Meméticos: AM-All vs AM-Rand vs AM-Best}

\begin{center}
\begin{tabular}{lrr|rr|rr}
\toprule
 & \multicolumn{2}{c|}{\textbf{ibex\_35}} & \multicolumn{2}{c|}{\textbf{syp\_100}} & \multicolumn{2}{c}{\textbf{syp\_500}} \\
\textbf{AM} & Fit.Train & Evals & Fit.Train & Evals & Fit.Train & Evals \\
\midrule
AM-All  & $-4.8652$ & $10000.0$ & $-3.4100$ & $10000.0$ & $-3.4839$ & $10000.0$ \\
AM-Rand & $-4.5161$ & $10000.2$ & $-3.1122$ & $10000.2$ & $-2.9504$ & $10000.2$ \\
AM-Best & $-4.5026$ & $10000.4$ & $-3.0789$ & $10000.6$ & $-2.9334$ & $10000.6$ \\
\bottomrule
\end{tabular}
\end{center}

\textbf{Hallazgo sorprendente}: \textbf{AM-All es peor que el mejor AG en todos los datasets}.
Ibex 35: AM-All ($-4.8652$) vs AGE-BLX ($-4.5205$), diferencia: $-0.345$.

\textbf{Causa}: AM-All aplica BL a $\approx 90$ individuos en la primera fase memetica
(generación 10), agotando prácticamente todo el presupuesto de 10000 evaluaciones.
El componente genético apenas evoluciona (solo 10 generaciones con $\approx 1000$ evals).
Los individuos todavía inmaduros quedan atrapados en óptimos locales por la BL prematura.
El resultado es peor que evolucionar solo genéticamente.

\textbf{AM-Best es superior}: Con solo 10 individuos por trigger, la BL es más selectiva.
Los mejores individuos se refinan a lo largo de múltiples fases meméticas intercaladas
con evolución genética. El componente genético tiene tiempo para diversificar.

\textbf{AM-Rand vs AM-Best}: AM-Rand selecciona aleatoriamente $\approx 10$ individuos.
AM-Best selecciona los mejores 10. AM-Best obtiene resultados ligeramente superiores
en todos los datasets (ibex\_35: $-4.5026$ vs $-4.5161$, diferencia $+0.013$),
confirmando que focalizar la explotación local en las mejores soluciones es más eficiente.

\subsection{Búsqueda Local vs Algoritmos Genéticos}

\textbf{La BL domina en todas las dimensiones}:

\begin{center}
\footnotesize
\begin{tabular}{lrrrr}
\toprule
\textbf{Dataset} & \textbf{BL} & \textbf{Mejor AG} & \textbf{Mejor AM} & \textbf{$\Delta$(BL$-$AG)} \\
\midrule
ibex\_35 & $-4.4912$ & AGE-BLX ($-4.5205$) & AM-Best ($-4.5026$) & $+0.029$ \\
syp\_100 & $-2.7911$ & AGE-BLX ($-2.9684$) & AM-Best ($-3.0789$) & $+0.177$ \\
syp\_500 & $-2.6282$ & AGG-BLX ($-2.8710$) & AM-Best ($-2.9334$) & $+0.243$ \\
\bottomrule
\end{tabular}
\end{center}

El gap crece dramáticamente con el tamaño del problema. En syp\_500, BL es $0.24$ puntos
mejor que el mejor AG, diferencia muy significativa.

\textbf{Razones}:
\begin{enumerate}
  \item \textbf{Operador problem-tailored}: La transferencia del 15\% de peso entre
  activos respeta las restricciones del portfolio y navega eficientemente el espacio
  factible. Los operadores genéricos (cruce/mutación) son menos especializados.

  \item \textbf{Presupuesto evaluativo}: Con 10000 evals y 100 individuos, el AG
  tiene $\approx 50$ generaciones de evolución. La BL usa todos los 10000 evals
  en una única trayectoria de mejora monotónica.

  \item \textbf{Convergencia de la BL}: En ibex\_35, la BL converge al óptimo local
  en tan solo $\approx 2100$ evals. En syp\_100 y syp\_500, agota las 10000 evals
  sin converger, lo que indica que el espacio es rico en mejoras locales accesibles.
\end{enumerate}

\subsection{Capacidad de generalización}

\textbf{Fitness de test (año 2025)}:

\begin{center}
\begin{tabular}{lcc}
\toprule
\textbf{Dataset} & \textbf{BL (test)} & \textbf{Mejor AG (test)} \\
\midrule
ibex\_35 & $-3.0850$ & AGE-BLX ($-3.1024$) \\
syp\_100 & $-3.7565$ & AGE-BLX ($-3.7770$) \\
syp\_500 & $-3.9162$ & AGG-BLX ($-3.9238$) \\
\bottomrule
\end{tabular}
\end{center}

\textbf{Observación}: La BL no solo obtiene mejor fit.train sino también mejor fit.test,
lo que indica que produce soluciones más \textbf{robustas} con mejor \textbf{generalización}.

Posible explicación: la trayectoria de mejora de la BL, al aplicar movimientos pequeños
y suaves (15\% de peso), tiende a converger hacia soluciones que respetan las correlaciones
entre activos, mientras que los cruzamientos y mutaciones genéticos pueden introducir
discontinuidades que sobreajustan al periodo de entrenamiento.

\subsection{Variabilidad y coste computacional}

\begin{center}
\begin{tabular}{lccc}
\toprule
\textbf{Algoritmo} & \textbf{ibex\_35} & \textbf{syp\_100} & \textbf{syp\_500} \\
\midrule
Random  & $0.024$s & $0.064$s & $1.734$s \\
BL      & $0.001$s & $0.039$s & $1.168$s \\
AGG-BLX & $0.008$s & $0.037$s & $0.673$s \\
AGE-BLX & $0.008$s & $0.034$s & $0.663$s \\
AM-Best & $0.016$s & $0.133$s & $1.326$s \\
\bottomrule
\end{tabular}
\end{center}

\textbf{Análisis de tiempos}:
\begin{itemize}
  \item En ibex\_35 (N=30), todos los algoritmos son rápidos ($\leq 0.025$s).
  La evaluación del fitness es O(N²) $\approx 900$ operaciones.

  \item En syp\_100 (N=97), la fitness es O(9409). Los AGs son ligeramente más rápidos
  que BL (0.034--0.037s vs 0.039s), ya que el operador de movimiento de la BL incurre
  en más overhead por evaluación que el cruce genético.

  \item En syp\_500 (N=457), la fitness es O(208849). AM-All es el más lento (2.111s)
  por la BL intensiva aplicada a $\approx 90$ individuos. Random (1.734s) es lento
  al generar 10000 soluciones completas desde cero. BL (1.168s) aplica movimientos
  incrementales baratos pero evalúa 10000 veces. Los AGs puros (0.654--0.673s) son
  los más rápidos al no incurrir en el overhead de la BL.
\end{itemize}

\textbf{Variabilidad}:
Los boxplots (Figuras \ref{fig:box_ibex}--\ref{fig:box_syp500}) muestran que:
\begin{itemize}
  \item \textbf{Random}: Alta varianza (muy disperso).
  \item \textbf{Greedy}: Sin varianza (determinístico).
  \item \textbf{BL}: Varianza moderada a baja.
  \item \textbf{AGs/AMs}: Varianza baja (el promediado de población reduce el ruido).
\end{itemize}

Los AGs tienden a ser más consistentes (menor varianza) que BL gracias al efecto de
averaging de la población, pero con peor valor medio.

\subsection{Conclusiones generales}

\begin{enumerate}
  \item \textbf{BLX supera a Aritmético} en ambos esquemas. La exploración fuera del
  convex hull es crítica para todos los tamaños de problema.

  \item \textbf{AGE supera a AGG} en presupuestos limitados. La conservación de élite
  implícita es más eficiente que el reemplazo generacional.

  \item \textbf{AM-Best y AM-Rand son competitivos pero no dominantes}. AM-All fracasa
  al concentrar todo el presupuesto de BL en la primera fase memetica con población inmadura.

  \item \textbf{BL domina globalmente}. El operador problem-tailored y la convergencia
  monotónica son decisivos con presupuestos limitados.

  \item \textbf{Escalabilidad}: El gap BL--AG crece con N. En problemas grandes,
  los métodos especializados ganan con claridad.
\end{enumerate}

%-----------------------------------------------------------------------
\end{document}
